% =====================================================================
% APPENDIX A (PRDL, Maria Martinez): use of AI scientific assistants
% Sources: docs/AI_USAGE_LOG.md (sessions, Bounty candidates B1, B13, B14,
% section 3), docs/literature/quotes.md (sections E and G), DECISIONS D8, D10.
% =====================================================================
\appendixbanner{APPENDIX A: USE OF AI SCIENTIFIC ASSISTANTS}
\addcontentsline{toc}{section}{Appendix A: Use of AI scientific assistants}\label{app:ai}

\subsection*{Workflow}

The repository file \texttt{docs/AI\_USAGE\_LOG.md} records each session with an assistant: the date and
conversation, who of us was present, the request, the output and the check we made. For literature the
tools never verified each other; each had one job:
\begin{enumerate}[leftmargin=1.4em,itemsep=1pt]
  \item \textbf{Perplexity} only listed candidate papers per topic (mine: recurrent and attention models,
        task T17), with no quotes and no numbers; we downloaded the PDFs ourselves.
  \item \textbf{NotebookLM}, limited to those PDFs, had to return for each claim a verbatim quote, the file
        and the printed page, or ``NOT FOUND''.
  \item \textbf{Claude} searched every quote in the text of its PDF, page by page; a quote enters the
        reports only from \texttt{docs/literature/quotes.md}, where it is marked as checked. Course
        materials support statements as course materials, never in place of the paper they mention.
  \item \textbf{Claude} then drafted the report from the block notes, the decisions and the checked quotes,
        and I went through each sentence against the logs, tables and figures; what no number supports is
        written as a hypothesis.
\end{enumerate}

\subsection*{Prompt strategy}

A Perplexity prompt asked for sources and forbade content, so that nothing it wrote could reach the report
by itself. An excerpt of the prompt for my topic (\texttt{docs/literature/prompts.md}, P2):
\begin{quote}\small\itshape
Gradient clipping for recurrent networks (Pascanu et al. 2013; Mikolov 2010), LSTM variants (Greff et al.
2015), GRU (Cho et al. 2014). Additive attention over RNN states (Bahdanau et al. 2015), the Transformer
(Vaswani et al. 2017), and Pre-LN versus Post-LN (Xiong et al. 2020). [\ldots] Mark ``no source found''
where that is the case. Do not give quotes or numbers. Do not invent identifiers.
\end{quote}
Its answer was a list of papers; Mikolov 2010 and Xiong et al.\ 2020 did not reach our PDF folder, which is why
those two claims were rewritten (below).

\subsection*{Claims of my blocks that were checked before writing}

While the recurrent and attention blocks were designed, the assistant attributed several statements to
primary papers that it had seen cited in the course slides, and flagged them itself as unverified. In the
literature stage each one was closed: the comparison of LSTM variants, the alignment model of additive
attention, gradient clipping and every component of the Transformer (scaled dot-product attention,
multi-head attention, Post-LN, sinusoidal encoding, dropout rate) now rest on checked quotes of the
original papers. Two claims had no PDF in the folder and were rewritten: gradient clipping is cited from
Pascanu et al.\ instead of ``first published'' in an earlier toolkit, and the easier optimisation of Pre-LN
is cited from the course slides that state it. The original LSTM paper and the papers behind the CLS token,
patch tokens and GELU were not in the folder either, so these components are described without
attribution. Explanations that are not measured (why the GRU beats the LSTM, why raw samples beat frames,
why positions do not help, the alignment of attention with the QRS) remain hypotheses.

\subsection*{Errors detected (candidates for the AI Error Bounty)}

Every assistant error of the project is documented in exactly one of our four reports. The three below
concern the data and the literature of my blocks:
\begin{itemize}[leftmargin=1.2em,itemsep=2pt]
  \item \textbf{A labelling rule stated from memory as the benchmark's} (Claude; B1; Sheet P1), in my data
        block, before any code.
  \item \textbf{Wrong page numbers in most answers} (NotebookLM; B13; Sheet P2), including quotes of my
        topic.
  \item \textbf{The article's licence given as the dataset's} (NotebookLM; B14; Sheet P3), a claim about the
        data of my block.
\end{itemize}
The other NotebookLM errors of the same checking stage (two versions of the benchmark paper mixed, a
frequency band offered for a duration, top-$k$ sampling attributed to a secondary source, statements
attributed to documents that were not uploaded) are documented in Oriana's PRDL report. The three errors
below were caught by the assistant while checking its own claim or another tool's answer, in checks that
we required: a factual claim needs a source (B1), and every NotebookLM quote is compared with its PDF
(B13, B14, decision D8).

\subsection*{AI Error Bounty sheets}

\subsubsection*{Sheet P1. A labelling rule stated from memory as the benchmark's (Claude)}
\begin{description}[style=nextline,leftmargin=1.2em]
  \item[Prompt] Which decisions were needed before writing the label code, and which option was best
        (conversation of the data block, 20 November).
  \item[Faulty response] Assign a superclass only for diagnostic statements with a likelihood above zero,
        presented as ``what the reference benchmark does'', without a source.
  \item[Detected by] The assistant itself, when we asked it to analyse the options before writing the
        label code.
  \item[Proof] The example script distributed with PTB-XL assigns the superclass whatever the likelihood;
        with that rule the class counts equal the table of the PhysioNet page, while the proposed rule loses
        STTC records and the class-count test fails.
  \item[Correction] Likelihood ignored (decision D2); a test of the suite compares the class counts of the
        cache with the published table.
  \item[Reproduction] \texttt{python -m src.prepare\_data} prints the counts. Changing the loop of
        \file{src/data.py::superclass_labels} so that it skips statements with likelihood 0 lowers the STTC
        count.
  \item[Reflection] ``This is what the benchmark does'' is a factual claim and needs the same check as a
        quote. Labelling rules are checked against the dataset's own files and published counts, because an
        error there changes every result of the project.
\end{description}

\subsubsection*{Sheet P2. Wrong page numbers (NotebookLM)}
\begin{description}[style=nextline,leftmargin=1.2em]
  \item[Prompt] Preamble of every NotebookLM prompt: for each claim, the literal quote, the file and ``the
        page number printed on the PDF page (not the viewer page)''.
  \item[Faulty response] The quotes were mostly correct, but most answers gave wrong pages, in several
        papers (the dataset paper, the dropout paper, CLIP, and in my topic the paper on the difficulty of
        training recurrent networks).
  \item[Detected by] The assistant, in the checking stage of decision D8.
  \item[Proof] Each quote was searched in the text of its PDF page by page, and the printed page was read
        on the page where it appears.
  \item[Correction] The pages in \texttt{docs/literature/quotes.md} are taken from the PDF, printed page
        first; no page given by NotebookLM is used.
  \item[Reproduction] Open each PDF at the page given in \texttt{quotes.md}; the quote is there.
  \item[Reflection] A page number looks like a detail, but without the right page nobody can check the
        quote; asking a tool for the page does not mean that it reads it.
\end{description}

\subsubsection*{Sheet P3. The article's licence given as the dataset's (NotebookLM)}
\begin{description}[style=nextline,leftmargin=1.2em]
  \item[Prompt] For the dataset paper, the licence of the data, with a literal quote, file and page.
  \item[Faulty response] A quote with the licence of the published article, given as the licence of the
        dataset.
  \item[Detected by] The assistant, in the checking stage of decision D8.
  \item[Proof] The quote is in the article's own copyright notice and says nothing about the data files.
  \item[Correction] Row not used (\texttt{quotes.md}, A16). The dataset licence had already been confirmed
        on the PhysioNet page by Oriana and me, and the report cites that page.
  \item[Reproduction] Read the licence notice at the end of the dataset paper and the licence on the
        PhysioNet page of PTB-XL 1.0.3.
  \item[Reflection] A quote can be literal and about the wrong object: the licence of an article is not the
        licence of the data it describes.
\end{description}

\subsection*{Reflection}

The assistants were most useful where their output can be checked mechanically: code with tests,
metrics recomputed from saved predictions, quotes searched in the PDF. They were least reliable when they
stated facts from memory or summarised a source. For the defence I studied each technique of my blocks
until I could explain it without the assistant, and I re-read the code line by line before running it.
